% !PRINTED 634
\weberpage{634}
%<<CONT
daher durch einige dieser Primfactoren theilbar sein, und sie können, so wenig wie $p$ selbst, durch das Quadrat eines $\mathfrak{p}_n$ theilbar sein. Wir müssen feststellen, durch welche Primtheiler $\mathfrak{p}_n$ jede dieser Zahlen $\psi_s(r)$ theilbar ist.

Dazu aber führen einerseits die Congruenzen des vorigen Paragraphen:
\weq{12}{r \equiv g^{-n\frac{p-1}{m}} \ (\mathrm{mod}\ \mathfrak{p}_{n'}), \qquad n n' \equiv 1\ (\mathrm{mod}\ m),}
andererseits die im \S.\,170 des ersten Bandes abgeleiteten Congruenzen (17):
\weq{13}{\begin{aligned}
&\psi_{\lambda,\mu}(g) \equiv 0, && \lambda+\mu<p-1\\[4pt]
&\psi_{\lambda,\mu}(g) \equiv -\dfrac{\Pi(2p-\lambda-\mu-2)}{\Pi(p-\lambda-1)\,\Pi(p-\mu-1)}, && (\mathrm{mod}\ p)\\
&&& p-1<\lambda+\mu<2p-2,
\end{aligned}}
worin $\lambda$ und $\mu$ zwischen $0$ und $p-1$ genommen sind, so dass der Binominalcoëfficient % [sic] "Binominalcoëfficient" as printed; "Binomialcoëfficient" is meant
\[
\frac{\Pi(2p-\lambda-\mu-2)}{\Pi(p-\lambda-1)\,\Pi(p-\mu-1)}
\]
eine durch $p$ nicht theilbare ganze Zahl ist.

Nach der Congruenz (12) ist
\[
\psi_s(r) \equiv \psi_s(g^{-nm'}) \equiv \sum\limits^{t} g^{-nm'\,[\mathrm{ind}\,t - (s+1)\,\mathrm{ind}\,(t+1)]},
\]
oder nach (6):
\[
\psi_s(r) \equiv \psi_{\lambda,\mu}(g) \ (\mathrm{mod}\ \mathfrak{p}_{n'}),
\]
wenn
\weq{14}{\begin{aligned}
\lambda &\ \text{den kleinsten positiven Rest von } {-}nm's\\
\mu &\ \text{\textquotedbl\ \textquotedbl\ \textquotedbl\ \textquotedbl\ \textquotedbl}\ {-}nm'
\end{aligned}\quad (\mathrm{mod}\ p-1)}
bedeutet, und die Congruenzen (13) zeigen dann, dass $\psi_s(r)$ durch $\mathfrak{p}_{n'}$ theilbar ist, wenn $\lambda+\mu<p-1$, und nicht theilbar, wenn $\lambda+\mu>p-1$ ist.

Da die zu $m$ theilerfremde Zahl $n$ nur nach dem Modul $m$ bestimmt ist, so nehmen wir sie jetzt \gesperrt{positiv und kleiner als $m$ an}. Dann ist
\[
\mu = m'(m-n).
\]

Die aus (14) bestimmte Zahl $\lambda$ ist ein Vielfaches von $m'$ und wenn wir $\lambda = m'\varkappa$ setzen, so ist
\weq{15}{\varkappa \ \text{der kleinste positive Rest von } {-}ns \ (\mathrm{mod}\ m).}

Hiernach wird $\psi_s(r)$ durch $\mathfrak{p}_{\varkappa'}$ theilbar sein, wenn
% [sic] Source prints the index varkappa' in this sentence; the preceding derivation and following factor count indicate mathfrak p_{n'}. Preserve the printed index.
\[
m'(\varkappa+m-n) < m m'
\]
oder $\varkappa < n$, und nicht theilbar, wenn $\varkappa > n$ ist.
